\chapter{Virtual memory}
\label{ch:vm}

Every process believes it owns a huge, private address space starting at
0. In reality its pages are scattered over physical memory, many of them
are not there at all, and some are shared with other processes. The MMU
translates every single memory access, and the kernel keeps the
translation tables. This chapter covers the translation (x86-64's
four-level page tables, the TLB), the page fault as the kernel's hook
into it (demand paging, copy-on-write), and what happens when physical
memory runs out (page replacement).

\section{Why virtual memory}

\begin{itemize}
  \item \textbf{Isolation}: a process can only reach the frames its page
    tables map; the kernel's pages are mapped without the user bit.
  \item \textbf{Relocation}: every program can be linked for the same
    addresses.
  \item \textbf{Overcommit}: memory is promised (\code{mmap} of 1~GiB) and
    delivered page by page on first touch; unused parts cost nothing.
  \item \textbf{Sharing}: the same physical frame can appear in several
    address spaces (shared libraries, copy-on-write after \code{fork}).
\end{itemize}

\section{Paging and the x86-64 page tables}

Memory is divided into 4~KiB \emph{pages} (virtual) and \emph{frames}
(physical). A page table maps page numbers to frame numbers. A flat table
for a 48-bit address space would have $2^{36}$ entries of 8 bytes ---
512~GiB per process. Instead, x86-64 uses a tree of four levels; each
table is one 4~KiB frame with 512 entries, and only the parts of the tree
that map something exist.

\begin{figure}[htbp]
\centering
\begin{tikzpicture}[font=\small,>=Latex,
    f/.style={draw,minimum height=0.6cm,align=center},
    tab/.style={draw,minimum width=1.6cm,minimum height=1.7cm,fill=black!4,align=center}]
  % virtual address fields
  \node[f,minimum width=1.9cm,fill=newbg] (i4) {PML4 idx\\[-2pt]\scriptsize 47--39};
  \node[f,minimum width=1.9cm,fill=newbg,right=0pt of i4] (i3) {PDPT idx\\[-2pt]\scriptsize 38--30};
  \node[f,minimum width=1.9cm,fill=newbg,right=0pt of i3] (i2) {PD idx\\[-2pt]\scriptsize 29--21};
  \node[f,minimum width=1.9cm,fill=newbg,right=0pt of i2] (i1) {PT idx\\[-2pt]\scriptsize 20--12};
  \node[f,minimum width=2.4cm,fill=keybg,right=0pt of i1] (off) {offset\\[-2pt]\scriptsize 11--0};
  \node[above=2pt of i3.north east] {virtual address};
  % tables
  \node[tab,below=1.2cm of i4] (t4) {PML4};
  \node[tab,below=1.2cm of i3] (t3) {PDPT};
  \node[tab,below=1.2cm of i2] (t2) {PD};
  \node[tab,below=1.2cm of i1] (t1) {PT};
  \node[tab,below=1.2cm of off,fill=keybg] (fr) {frame\\(4 KiB)};
  \node[left=0.7cm of t4,draw,fill=warnbg] (cr3) {CR3};
  \draw[->] (cr3) -- (t4);
  \foreach \i/\t in {i4/t4,i3/t3,i2/t2,i1/t1,off/fr} \draw[->,dashed] (\i) -- (\t);
  \draw[->,thick] (t4) -- (t3); \draw[->,thick] (t3) -- (t2);
  \draw[->,thick] (t2) -- (t1); \draw[->,thick] (t1) -- (fr);
\end{tikzpicture}
\caption{A page walk: each 9-bit index selects one of 512 entries; the
entry holds the frame number of the next table, and finally of the page.}
\label{fig:pagewalk}
\end{figure}

A page-table entry (64 bits) holds the frame number (bits 12--51) and
flags, among them:
\begin{center}
\begin{tabular}{lll}
\toprule
\textbf{Bit} & \textbf{Name} & \textbf{Meaning} \\
\midrule
0 & P (present) & the entry is valid; otherwise any access faults \\
1 & W (writable) & stores allowed \\
2 & U (user) & accessible from ring 3 \\
5 & A (accessed) & set by the MMU on every access \\
6 & D (dirty) & set by the MMU on every write \\
9--11 & available & ignored by the MMU, free for the OS (e.g.\ a COW mark) \\
63 & NX & no instruction fetch \\
\bottomrule
\end{tabular}
\end{center}

The effective rights are the AND over all four levels: the upper levels
are usually created permissive (P, W, U) and the leaf entry decides.
Mapping one page into an empty address space costs three new tables
(PDPT, PD, PT); unmapping should free tables that became empty.

\section{The TLB}

A page walk costs four extra memory accesses. The \emph{translation
lookaside buffer} caches recent translations (page number $\to$ frame,
rights) and makes almost all accesses hit. The TLB is not coherent with
the page tables: when the kernel changes or removes a present entry, it
must invalidate the cached copy (\code{invlpg addr}, or reload \code{cr3}
to drop everything). Forgetting it leaves the old translation usable ---
after an unmap, a program keeps reading a frame that may already belong to
someone else. Entries that were not present are never cached, so
\emph{adding} a mapping needs no invalidation. On multi-core machines
other CPUs may cache the translation too: a \emph{TLB shootdown} sends
them an interrupt to invalidate --- expensive, and batched by real
kernels.

\section{Page faults}

When the MMU cannot translate (not present) or the rights do not allow
the access (write to read-only, user access to kernel page), it raises a
page fault: the CPU stores the faulting address (x86: in \code{cr2}) and
an error code (present?, write?, user?) and enters the kernel's handler.
The handler decides:

\begin{itemize}
  \item legal but not loaded yet: allocate/load the page, map it, return
    --- the CPU \emph{re-executes} the faulting instruction;
  \item copy-on-write page written: copy it, map the copy writable, return;
  \item anything else: the program made a mistake --- signal
    \code{SIGSEGV} (Linux) or end the process (SWEB).
\end{itemize}

\paragraph{Demand paging.} Nothing is loaded before it is used. Code
and data are read from the executable on the first access; heap and stack
pages are allocated on first touch and \textbf{zero-filled} --- a frame
straight from the free pool may contain another process's data. On Linux
even reading an untouched anonymous page only maps a shared zero page.

\paragraph{Copy-on-write.} \code{fork} shares all frames between parent
and child and makes every writable page read-only in \emph{both}, with a
software ``COW'' mark. The first write from either side faults; the
handler copies the frame and maps the copy writable. If the frame has only
one user left (reference count 1), it simply becomes writable again. Note
that the parent's TLB may still hold writable entries for those pages ---
they must be invalidated, or the parent's writes land in the shared frame.

\section{Page replacement}

When no frame is free, a page must be evicted (written to swap if dirty)
to make room. Which one?

\begin{description}
  \item[OPT (Belady)] evict the page whose next use lies furthest in the
    future. Optimal, but needs the future; the yardstick for the others.
  \item[FIFO] evict the page loaded longest ago. Simple, often bad, and it
    suffers from \emph{Belady's anomaly}: more frames can cause more
    faults.
  \item[LRU] evict the page unused for the longest time. Good, but exact
    LRU needs a timestamp per access --- too expensive in hardware.
  \item[Clock (second chance)] approximates LRU with the accessed bit: a
    hand sweeps over the frames; a page with R = 1 gets its bit cleared
    and a second chance, the first page with R = 0 is evicted.
\end{description}

Worked example, reference string \code{7 0 1 2 0 3 0 4 2 3 0 3 2 1 2 0 1 7
0 1}, 3 frames: FIFO 15 faults, LRU 12, OPT 9, Clock 14. Belady's anomaly
with FIFO on \code{1 2 3 4 1 2 5 1 2 3 4 5}: 9 faults with 3 frames, 10
with 4. LRU and OPT are \emph{stack algorithms} --- the pages kept with $k$
frames are always a subset of those kept with $k+1$ --- and never show the
anomaly. (\module{08}, part C simulates all four with exact rules.)

\paragraph{Working set and thrashing.} The \emph{working set} is the set
of pages a process used in the recent past. If the working sets of all
running processes do not fit into memory, every process keeps faulting
out the pages the others need: the system is busy paging and does no
useful work --- \emph{thrashing}. The remedy is fewer processes in memory
(suspend some), not a faster disk.

\begin{swebbox}
\file{arch/x86/64/source/ArchMemory.cpp}: \code{mapPage} walks the four
levels and allocates missing tables (P, W, U on the upper levels);
\code{unmapPage} frees tables that became empty; frames come from
\code{PageManager::allocPPN}. The kernel reaches any frame through the
identity mapping at \code{IDENT\_MAPPING\_START}
(\code{getIdentAddressOfPPN}). Page faults go to
\file{common/source/mm/PageFaultHandler.cpp}; a valid user fault calls
\code{Loader::loadPage}, which only knows the ELF segments --- a fault
elsewhere, e.g.\ just below the single stack page, ends the process.
For P1 this matters twice: every new thread needs its own user stack
somewhere in the shared address space, and two threads of one process can
now fault at the same time --- \code{mapPage} has no lock of its own.
\end{swebbox}

\begin{keybox}
\begin{itemize}
  \item Four levels of 512-entry tables; the leaf decides the rights, the
    MMU ANDs all levels.
  \item Changing or removing a present PTE requires a TLB invalidation.
  \item Page faults are the kernel's hook: demand paging (zero-filled!),
    copy-on-write, or kill.
  \item COW: read-only + mark in both address spaces, copy on first write,
    no copy for the last user.
  \item OPT $\le$ LRU; FIFO shows Belady's anomaly; Clock approximates LRU
    with the accessed bit.
  \item Thrashing: working sets do not fit; reduce the load.
\end{itemize}
\end{keybox}
